\begin{thm} Let $p$ be an odd prime. Then ${-1 \choose p} = (-1)^{(p-1)/2}$. \end{thm}
\begin{proof}
	Consider the ideal $p\ZZ[i] \subseteq \ZZ[i]$. Since $p \nmid 4 = \disc(\ZZ[i])$, we have that $p$ is unramified, so $p$ splits as a product
	\[ p\ZZ[i] = P_1\cdots P_r \]
	in $\ZZ[i]$. Then, the sum of the inertial degrees is 2 since we have a degree 2 extension. I.e. we either have $p\ZZ[i] = P_1P_2$ with $f(P_1|p) = f(P_2|p) = 1$ or else $P_1 = p\ZZ[i]$ is itself prime of inertial degree 2. In either case, we compute the inertial degree directly.
	
	Namely, let $P$ be a prime lying over $p$, and let $f$ be the inertial degree. Then, by definition, $\ZZ[i]/P$ is a degree $f$ extension of $\ZZ/p\ZZ = \FF_p$. We thus know that the Galois group is generated by the Frobenius automorphism $x \mapsto x^p$. But any map of $\ZZ[i]/P$ is determined by the image of $i$, and $i^p = \pm i$ with the plus sign holding iff $p \equiv 1 \pmod{4}$. Thus, $f$, the order of the Galois group, is 1 iff $p \equiv 1 \pmod{4}$.
	
	That is, we have $p \equiv 1 \pmod{4}$ iff $p\ZZ[i]$ is not prime. But we can characterize this as well:
	\[ \ZZ[i]/p\ZZ[i] \cong \ZZ[x]/(x^2+1,p) \cong \FF_p/(x^2+1) \]
	so $p\ZZ[i]$ is not prime iff $x^2+1$ is reducible in $\FF_p$. Since it is of degree 2, this can only happen by having a root, i.e. $p \equiv 1 \pmod{4}$ iff there is some $a \in \FF_p$ with $a^2+1=0$, i.e. $a^2=-1$, so $-1$ is a quadratic residue. This gives the result.
\end{proof}

\begin{thm} Let $p,q$ be distinct odd primes. Then ${p \choose q}{q \choose p} = (-1)^{(p-1)(q-1)/4}$. \end{thm}
\begin{proof}
	If $q \equiv 3 \pmod{4}$, let $m = q$, else let $m = -q$. Then $m\ZZ = q\ZZ$ and $m \equiv 3 \pmod{4}$, and in particular, the ring of integers in $\QQ(\sqrt{m})$ is $R = \ZZ[\sqrt{m}]$. We mimic the above proof in this case.
	
	The ideal $pR$ is unramified since $p \nmid \disc(R) = 4m = \pm 4q$. So, as above, we have two cases since the extension is of degree 2: there are either one or two primes of $R$ lying over $p$, and there is one prime iff $pR$ is itself prime of inertial degree $2$.
	
	But then again, if $P$ is a prime of $R$ lying over $p$, then $R/P$ is a degree $f$ extension of $\FF_p$ and so is generated by Frobenius. 
\end{proof}
